\section{Síntesis y ampliación}

\subsection{Puntos clave de esta clase}

Esta clase es la introducción al curso KAIST CS492D, estableciendo sistemáticamente el marco teórico para los modelos generativos:

\begin{enumerate}
    \item \textbf{Problema central de los modelos generativos}: Dado un conjunto de muestras de datos, aprender la distribución subyacente $p_{\text{data}}(x)$ para generar nuevas muestras.
    \item \textbf{Marco unificado}: Todos los modelos de generación profunda aprenden una mapeo desde una distribución simple $p(z)$ hacia la distribución de datos $p_{\text{data}}(x)$.
    \item \textbf{GAN}: Aprende el mapeo mediante un juego entre generador y discriminador (optimización minimax). Teóricamente elegante, pero su entrenamiento es extremadamente inestable.
    \item \textbf{VAE}: Aprende la distribución condicional $p_\theta(x|z)$ mediante inferencia variacional, maximizando el ELBO. Su entrenamiento es estable y sienta las bases para los Diffusion Models.
    \item \textbf{Herramientas matemáticas}: Muestreo por transformación inversa, transformación Box-Muller, truco de reparametrización, marginalización, teorema de Bayes, divergencia KL, desigualdad de Jensen.
\end{enumerate}

\subsection{Comparación GAN vs. VAE}

\begin{table}[H]
\centering
\begin{tabular}{lll}
\toprule
Característica & GAN & VAE \\
\midrule
Método de entrenamiento & Entrenamiento adversario Minimax & Maximización del ELBO \\
Estabilidad del entrenamiento & Inestable, propenso a fallos & Relativamente estable \\
Calidad generada & Alta (especialmente nitidez en imágenes) & Media (tiende a la difuminación) \\
Diversidad & Propenso al colapso de modos & Mejor \\
Fundamento teórico & Teoría de juegos & Inferencia variacional \\
¿Posee encoder? & No (solo generador + discriminador) & Sí (encoder + decoder) \\
Estructura del espacio latente & Sin restricciones explícitas & Con restricción regularizadora KL \\
Relación con Diffusion & Débil & Fuerte (el marco ELBO se generaliza directamente) \\
\bottomrule
\end{tabular}
\caption{Comparación sistemática entre GAN y VAE}
\end{table}

\subsection{Perspectivas: Modelos de Difusión y Flujo}

Esta clase prepara el terreno con toda la base teórica necesaria para los contenidos posteriores. En las siguientes clases, el curso explorará en profundidad:

\begin{itemize}
    \item \textbf{Diffusion Model}: Puede entenderse como un ``VAE de múltiples pasos'', que modela la distribución de datos mediante un proceso gradual de adición y eliminación de ruido. La derivación del ELBO se puede generalizar directamente al contexto de difusión.
    \item \textbf{Flow Model}: Modela la verosimilitud logarítmica exacta mediante transformaciones invertibles (sin necesidad de una aproximación inferior como el ELBO).
    \item \textbf{Score-based Model}: Logra la generación desde otra perspectiva, aprendiendo la función score de la distribución de datos (el gradiente de la densidad logarítmica).
\end{itemize}

Los conceptos teóricos de muestreo de esta clase (muestreo por transformación inversa, truco de reparametrización), el marco probabilístico (prior / verosimilitud / posterior) y la derivación del ELBO aparecerán y se profundizarán en cada una de las siguientes clases.

\subsection{Lecturas adicionales}

\begin{itemize}
    \item Goodfellow et al., \textit{Generative Adversarial Nets}, NeurIPS 2014
    \item Kingma \& Welling, \textit{Auto-Encoding Variational Bayes}, ICLR 2014
    \item Kingma \& Welling, \textit{An Introduction to Variational Autoencoders}, Foundations and Trends in ML, 2019
    \item Doersch, \textit{Tutorial on Variational Autoencoders}, arXiv:1606.05908
    \item Página principal del curso KAIST CS492D: \url{https://diffusion.kaist.ac.kr/}
\end{itemize}

%% --- Fin del contenido --- %%

\end{document}
